% Source: https://artofproblemsolving.com/community/c3h2544460_trig_help
% Stable identity: aops-topic-2544460-q01
\textbf{(a)} Reflecting in the $y$-axis maps the unit circle to itself and sends $(x,y)$ to $(-x,y)$. It maps the positive $x$-axis to the negative $x$-axis and reverses the sense of turning, so the ray at angle $\theta$ goes to the ray making angle $\theta$ with the negative $x$-axis, measured clockwise; that is the ray at angle $180^\circ-\theta$. (In the sketch, the red radius makes the angle $\theta$ with the negative $x$-axis.) Hence the point $(\cos\theta,\sin\theta)$ goes to $\left(\cos\left(180^\circ-\theta\right),\sin\left(180^\circ-\theta\right)\right)$, and this equals $(-\cos\theta,\sin\theta)$. $\blacksquare$

\textbf{(b)} In the same way, a half-turn about the origin gives $(\cos(180^\circ+\theta),\sin(180^\circ+\theta))=(-\cos\theta,-\sin\theta)$, and reflection in the $x$-axis gives $(\cos(360^\circ-\theta),\sin(360^\circ-\theta))=(\cos\theta,-\sin\theta)$. Using the values for $30^\circ$, $45^\circ$, $60^\circ$:
\[\sin150^\circ=\sin30^\circ=\boxed{\dfrac12},\qquad\cos150^\circ=-\cos30^\circ=\boxed{-\dfrac{\sqrt3}{2}},\]
\[\sin225^\circ=-\sin45^\circ=\boxed{-\dfrac{\sqrt2}{2}},\qquad\cos300^\circ=\cos60^\circ=\boxed{\dfrac12},\]
\[\tan120^\circ=\dfrac{\sin60^\circ}{-\cos60^\circ}=\dfrac{\dfrac{\sqrt3}{2}}{-\dfrac12}=\boxed{-\sqrt3}.\]

\begin{asy}
usepackage("amsmath");
size(5.5cm);
defaultpen(fontsize(10pt));
draw((-1.35,0)--(1.4,0),Arrow(TeXHead));draw((0,-1.3)--(0,1.35),Arrow(TeXHead));
draw(unitcircle,linewidth(0.8));
pair P=dir(30), Q=dir(150);
draw((0,0)--P,blue+linewidth(0.8));draw((0,0)--Q,red+linewidth(0.8));
draw(P--(P.x,0),blue+dashed);draw(Q--(Q.x,0),red+dashed);
draw(arc((0,0),0.25,0,30),blue);draw(arc((0,0),0.25,150,180),red);
dot(P,blue);dot(Q,red);
label("$(\cos\theta,\sin\theta)$",P,NE,blue+fontsize(8pt));
label("$(-\cos\theta,\sin\theta)$",Q,NW,red+fontsize(8pt));
label("$\theta$",0.25*dir(15),E,blue+fontsize(8pt));
label("$\theta$",0.25*dir(165),W,red+fontsize(8pt));
label("$x$",(1.4,0),E);label("$y$",(0,1.35),N);
\end{asy}
